% =====================================================================
\section{Cross-dataset fit}\label{sec:cross-dataset}
% =====================================================================

The kernel is now applied, with no shape change, to four independent
$B \to K^{*}\mu\mu$ datasets covering two collaborations, two
isospin partners, and two collision energies. Each dataset uses its
own $q^{2}$ bin grid, its own published correlation matrix where
available, and its own per-bin SM and non-linear new-physics
predictions regenerated by \texttt{flavio} 2.4.

\subsection{Datasets}

\begin{table}[h]
\centering
\small
\begin{tabular}{llllrl}
\toprule
dataset & decay & arXiv & HEPData & lumi (fb$^{-1}$) & source bins \\
\midrule
LHCb 2015         & $B^{0}\to K^{*0}\mu\mu$  & 1512.04442 & ins1409497 & 3   & 8 (Table 4 P-basis) \\
LHCb 2021         & $B^{+}\to K^{*+}\mu\mu$  & 2012.13241 & ins1838196 & 9   & 8 exclusive \\
CMS 2025          & $B^{0}\to K^{*0}\mu\mu$  & 2410.18247 & ins2850101 & 140 & 6 \\
LHCb 2025 (ref)   & $B^{0}\to K^{*0}\mu\mu$  & 2512.18053 & ins3094698 & 8.4 & 8 \\
\bottomrule
\end{tabular}
\caption{Datasets used for the cross-dataset test.}
\label{tab:cross_dataset_summary}
\end{table}

LHCb 2021 publishes 10 $q^{2}$ bins of which two are wide combined
($[1.1,6.0]$ and $[15,19]$); these overlap the exclusive ones and are
dropped. The covariance from HEPData is block-diagonal across $q^{2}$.
CMS 2025 publishes block-diagonal correlation for 6 $q^{2}$ bins
(charmonium windows excluded by their analysis).

\subsection{Headline result (non-linear)}\label{subsec:cross_headline_nl}

The four observables retained for the joint fit on each dataset are
$P_{5}', P_{4}', P_{1}, P_{2}$, matching the LHCb 2025 reference fit.
Same kernel function evaluated at each dataset's own bin centres. All
fits are evaluated using \texttt{flavio.np\_prediction} directly
(no Taylor expansion).

\begin{table}[h]
\centering
\footnotesize
\setlength{\tabcolsep}{3.5pt}
\begin{tabular}{lcrrrrrr}
\toprule
dataset & $n$ & FREE $\chi^{2}$ & FREE $\Delta C_{9}$ &
        VFD $\chi^{2}$ & \textbf{VFD $A$} & VFD $\Delta C_{9}^{\mathrm{eff}}$ &
        $\Delta\mathrm{AIC}$ \\
\midrule
LHCb 2015                     & 32 & 30.69 & $-1.08$ & 30.45 & $\boldsymbol{+1.24}$ & $-0.96$ & $-0.24$ \\
LHCb 2021 ($B^{+}\to K^{*+}$) & 32 & 22.77 & $-1.82$ & 22.93 & $\boldsymbol{+2.06}$ & $-1.59$ & $+0.17$ \\
CMS 2025 (no $P_{4}'$)        & 18 & 43.73 & $-0.95$ & 44.21 & $\boldsymbol{+1.05}$ & $-0.81$ & $+0.47$ \\
LHCb 2025 (ref)               & 32 & 40.89 & $-1.00$ & 41.98 & $\boldsymbol{+1.14}$ & $-0.86$ & $+1.09$ \\
\midrule
$B_{s}\!\to\!\phi$ 2015        & 24 & 13.20 & $-4.12$ & 12.96 & $\boldsymbol{+4.98}$ & $-3.85$ & $-0.24$ \\
\bottomrule
\end{tabular}
\caption{Headline non-linear cross-dataset and cross-channel results.
At equal $k=1$ the kernel is favoured on $2/5$ fits (LHCb 2015 and
$B_{s}\!\to\!\phi$, both by $\Delta\mathrm{AIC}=-0.24$) and disfavoured
on $3/5$ (LHCb 2021, CMS 2025 no-$P_{4}'$, LHCb 2025; max
$\Delta\mathrm{AIC}=+1.09$). $A>0$ in every fit;
$\Delta C_{9}^{\mathrm{eff}}<0$ in every fit. Reproducible from
\texttt{scripts/wo016d\_nonlinear\_xdataset.py}
(reports/wo016d\_nonlinear\_xdataset.csv).}
\label{tab:cross_dataset_main}
\end{table}

\begin{figure}[h]
\centering
\includegraphics[width=0.95\textwidth]{fig_F3_cross_dataset_A.pdf}
\caption{Non-linear best-fit amplitudes $A$ across the five fits.
Left: the four $P$-basis $B\to K^{*}\mu\mu$ fits, range
$A\in[+1.05,+2.06]$, factor-of-$\sim 2$ scatter across two
collaborations and two isospin partners. Right: the $S$-basis
$B_{s}\!\to\!\phi$ fit, $A=+4.98$. The grey dashed line marks the
$P\!\to\!S$ basis-corrected prediction
$A_{P}^{\mathrm{LHCb 2025}}\times 2.2 \approx 2.5$;
the observed $S$-basis amplitude exceeds this by a factor of $\sim 2$
(see \S\ref{sec:cross-channel}). Bars are coloured green where the
non-linear $\Delta\mathrm{AIC}$ marginally favours VFD and orange
where it marginally favours $\mathrm{FREE\_C9}$. Reproducible from
\texttt{scripts/wo017\_paper\_figures.py}.}
\label{fig:F3_cross_dataset_A}
\end{figure}

\subsection{Akaike-weight stack across the five fits}

To convert the per-dataset $\Delta\mathrm{AIC}$ values into a single
combined statement we compute the Akaike weight of each model on each
dataset, $w_{M,d} = e^{-\Delta\mathrm{AIC}_{M,d}/2} / \sum_{M'}
e^{-\Delta\mathrm{AIC}_{M',d}/2}$, and stack across datasets under the
independence assumption (none of the five datasets shares
observation-level information; theory inputs are shared via
\texttt{flavio} but enter both arms identically).

\begin{table}[h]
\centering
\small
\begin{tabular}{lcccc}
\toprule
dataset & non-linear $\Delta\mathrm{AIC}$ &
          $w(\mathrm{FREE\_C9})$ & $w(\mathrm{VFD})$ & favoured \\
\midrule
LHCb 2015                  & $-0.241$ & $0.470$ & $0.530$ & VFD (marginal) \\
LHCb 2021 ($K^{*+}$)       & $+0.168$ & $0.521$ & $0.479$ & FREE\_C9 (marginal) \\
CMS 2025 no-$P_{4}'$       & $+0.473$ & $0.559$ & $0.441$ & FREE\_C9 (marginal) \\
LHCb 2025                  & $+1.093$ & $0.633$ & $0.367$ & FREE\_C9 \\
$B_{s}\!\to\!\phi$ 2015    & $-0.240$ & $0.470$ & $0.530$ & VFD (marginal) \\
\midrule
\textbf{stacked}            & $\sum=+1.253$ &
\textbf{0.652} & \textbf{0.348} & \textbf{FREE\_C9 (mild)} \\
\bottomrule
\end{tabular}
\caption{Per-dataset Akaike weights and stacked total. The
multiplicative stack across five fits gives a $\sim 1.9{:}1$ odds in
favour of the constant $C_{9}$-shift. Reproducible from
\texttt{scripts/wo016a\_akaike\_stack.py}
(reports/wo016a\_akaike\_stack.md).}
\label{tab:akaike_stack}
\end{table}

\paragraph{Sign-uniformity test.} Across all five fits, $A>0$ and
$\Delta C_{9}^{\mathrm{eff}}<0$. Under the random-sign null hypothesis
$\mathbb{P}(A>0\,|\,\text{noise only}) = 1/2$, the probability of
$5/5$ positive amplitudes is $2^{-5}\approx 0.031$. \emph{The
appropriate null here is the ``no anomaly'' null}, against which
$\mathrm{FREE\_C9}$ would also produce $5/5$ same-sign fits on the
same data — the $b\to s$ angular anomaly is biased a priori to
$\Delta C_{9}<0$, and any sufficiently centre-peaked one-parameter
shape applied to it would inherit that bias. The sign-uniformity
calculation is therefore a consistency check on the kernel's
direction across collaborations and channels, \emph{not} independent
evidence for the kernel over a constant shift; the data establishes
the anomaly itself at much higher significance than $2^{-5}$ via
$\mathrm{FREE\_C9}$ alone.

\paragraph{Cross-dataset amplitude scatter.} The non-linear fitted
amplitudes on the four $P$-basis $B\to K^{*}$ fits span
$A\in[+1.05,+2.06]$ — a factor $\sim 2$ across two collaborations
and two isospin partners. Including the $S$-basis $B_{s}\!\to\!\phi$
fit ($A=+4.98$) widens the spread to a factor $\sim 5$, but
\S\ref{sec:cross-channel} shows that this last gap is mostly a
basis effect (Krüger--Matias $P_{i}$ vs CP-averaged $S_{i}$).

\subsection{Acceptance gates}

\begin{table}[h]
\centering
\small
\begin{tabular}{lc}
\toprule
criterion & result (non-linear) \\
\midrule
\textbf{$\mathrm{VFD\ AIC} \leq \mathrm{FREE\_C9\ AIC}$ on every dataset}
   & \textbf{\textsc{fail}} (passes on $2/5$; max $\Delta\mathrm{AIC}=+1.09$) \\
\textbf{stacked Akaike weight $w_{\mathrm{VFD}} \geq 0.5$}
   & \textbf{\textsc{fail}} ($w_{\mathrm{VFD}} = 0.348$) \\
\midrule
$A$ same sign across datasets & \textsc{pass} (5/5 positive) \\
$\Delta C_{9}^{\mathrm{eff}}$ same sign & \textsc{pass} (5/5 negative) \\
$A$ order $\sim 1$ across $P$-basis datasets & \textsc{pass} (range $[+1.05,+2.06]$, factor 2) \\
no per-dataset retuning & \textsc{pass} (same kernel function everywhere) \\
\bottomrule
\end{tabular}
\caption{Cross-dataset acceptance gates under the headline non-linear
analysis. The two strict AIC-preference gates (top) \textsc{fail}
under non-linear evaluation; the four direction-consistency and
parsimony gates (bottom) \textsc{pass}. The headline is therefore:
the kernel is sign-consistent with the anomaly but is not the
preferred model on AIC.}
\end{table}

\subsection{The CMS $P_{4}'$ convention question}

Including CMS 2025's published $P_{4}'$ values in the linearised
Mode-B fit produced $\chi^{2}=169$ over 24 data points — about
$7\times$ per data point. The residuals are dominated by $P_{4}'$
entries falling \emph{outside} the conventional $[-1, +1]$ physical
range, e.g.\ $P_{4}'(14.18, 16) = -1.16 \pm 0.06$. CMS's normalisation
appears to differ from \texttt{flavio}'s convention. We retain
the no-$P_{4}'$ row in Table~\ref{tab:cross_dataset_main} as the clean
comparison and treat the $P_{4}'$-included row only as a diagnostic of
the convention mismatch. Both $\Delta\mathrm{AIC}$ values agree
qualitatively (FREE\_C9 favoured by $\sim 0.5$ AIC unit either way)
under non-linear refit; the no-$P_{4}'$ result is the one quoted in
the headline.

\subsection{What this section establishes}

The same geometry-derived kernel, with one fitted amplitude per
dataset and no shape retuning, is sign-uniform across five
independent flavour-physics measurements covering two collaborations
(LHCb, CMS), two isospin partners ($B^{0}\!\to\!K^{*0}$,
$B^{+}\!\to\!K^{*+}$), two decay channels ($K^{*}$ and $\phi$), and
ten years of analysis ($2015$ -- $2025$). However, on a strict
AIC criterion at equal parameter cost, the kernel does not improve
over a constant Wilson-coefficient shift on the headline non-linear
analysis: the stacked Akaike weight is $\approx 1.9$:$1$ in favour
of $\mathrm{FREE\_C9}$. The substantive result is direction
consistency under random-sign null
($p\approx 0.03$), not statistical preference of the kernel shape
over the constant shift on this set of fits.
